\section{Method}
\label{sec:method}

\begin{figure}[t]
\centering
\includegraphics[width=\columnwidth]{figures/fig_task_scenario.png}
\caption{Task scenario. Before the event, ascending and descending SAR observations provide complementary views. During rapid post-event response, the target orbit is available while the complementary orbit may not yet be available.}
\label{fig:scenario}
\end{figure}

\begin{figure*}[!t]
\centering
\includegraphics[width=\textwidth]{figures/fig_architecture_grsl.png}
\caption{Detailed Teacher--Student search architecture from the original manuscript. The counter orbit guides the Teacher search policy, while retrieved values remain target-side. The auxiliary TGD branch shown in the diagram is optional and is not used for the principal geometry-free Student comparison.}
\label{fig:architecture}
\end{figure*}

\subsection{Problem Formulation}
For each paired training sample, $X_t$ and $X_c$ denote target- and counter-orbit pre/post-event SAR stacks and $Y$ is the landslide mask. The Counter-Guided Search Teacher receives $(X_t,X_c)$ during training; the Target-Only Search Student receives $X_t$ during training and inference. The role of $X_c$ is to teach an evidence-search operation executable from $X_t$. Geometry codes are used to define diagnostic evaluation regions. The primary Student does not receive geometry as an input or use it to weight its loss.

\subsection{Counter-Guided Target Evidence Search Teacher}
An encoder produces multi-scale target and counter features,
\begin{equation}
T=E(X_t), \qquad C=E(X_c).
\label{eq:teacher_features}
\end{equation}
At each reference location $p$ and scale, the Teacher policy head predicts $K$ offsets and normalized weights from both inputs,
\begin{equation}
\pi_T(p)=\{\Delta p^T_{p,k},a^T_{p,k}\}_{k=1}^{K}=\mathcal{P}(T,C)_p,
\quad \sum_k a^T_{p,k}=1.
\label{eq:teacher_policy}
\end{equation}
The search operator reads target features only:
\begin{equation}
R_T(p)=\mathcal{S}(T;\pi_T)_p=
\sum_{k=1}^{K}a^T_{p,k}W_vT(p+\Delta p^T_{p,k}),
\label{eq:teacher_search}
\end{equation}
where fractional locations use interpolation and $W_v$ is a learned value projection. A decoder maps $R_T$ to $P_T$. This is the central constraint: the counter orbit guides policy generation, but no counter feature enters the searched value path or decoder.

\subsection{Target-Only Student and Search Distillation}
The Student forms target features, policy, and searched representation as
\begin{equation}
S=E_S(X_t), \qquad \pi_S=\mathcal{P}_S(S), \qquad R_S=\mathcal{S}(S;\pi_S).
\label{eq:student_search}
\end{equation}
It receives two transfer signals. Policy distillation teaches \emph{where} to search by aligning teacher and student offsets and weights, $\widetilde{\mathcal{L}}_{\mathrm{policy}}(\pi_S,\pi_T)$. Evidence distillation teaches \emph{what} to retrieve by applying both policies to the \emph{same student target memory}:
\begin{equation}
R_{T\rightarrow S}=\mathcal{S}(S;\pi_T), \qquad
\widetilde{\mathcal{L}}_{\mathrm{evidence}}=\ell(R_S,R_{T\rightarrow S}).
\label{eq:shared_memory}
\end{equation}
Thus, the Student neither reconstructs counter-orbit features nor receives a hidden fusion path; it learns how the privileged Teacher policy should be realized on target-side information available at deployment.

\subsection{Learning Objectives}
The primary Teacher uses pixelwise binary cross-entropy (BCE) segmentation loss without geometry-specific weighting. The selected Student first learns from target-only segmentation and subsequently adds policy and evidence alignment. Its objective is
\begin{equation}
\mathcal{L}_S=\mathcal{L}_{\mathrm{BCE}}+
\lambda_p\mathcal{L}_{\mathrm{policy}}+
\lambda_e\mathcal{L}_{\mathrm{evidence}},
\label{eq:student_loss}
\end{equation}
where policy alignment uses the Jensen--Shannon distance between sampling distributions and evidence alignment uses cosine distance between representations retrieved from the same Student target memory. The weights are calibrated from training losses and decay during the second stage.
